% tab:results

\begin{table}[!htbp]
\centering\small\setlength{\tabcolsep}{2pt}
\begin{tabular}{@{}p{4.16cm}p{1.61cm}p{2.83cm}p{1.28cm}p{1.61cm}p{2.83cm}p{1.28cm}@{}}
\toprule
 & \multicolumn{3}{c}{$w_1$} & \multicolumn{3}{c}{$v_1$} \\
Вариант & Расчёт & ИМ $\pm$ ДИ & $\Delta$, \% & Расчёт & ИМ $\pm$ ДИ & $\Delta$, \% \\
\midrule
\multicolumn{7}{l}{Коэффициент вариации: 0,8} \\
RCS, удаление & 0,4104 & 0,4162 ± 0,0083 & −1,39 & 2,0850 & 2,0907 ± 0,0107 & −0,27 \\
Очищение & 0,1925 & 0,1932 ± 0,0022 & −0,38 & 1,4830 & 1,4835 ± 0,0028 & −0,03 \\
RCS, новая длительность & 1,2060 & 1,1999 ± 0,0181 & 0,51 & 3,4023 & 3,3957 ± 0,0203 & 0,20 \\
Все каналы, новая длительность & 1,8598 & 1,8522 ± 0,0289 & 0,41 & 4,1856 & 4,1759 ± 0,0298 & 0,23 \\
\multicolumn{7}{l}{Коэффициент вариации: 1,2} \\
RCS, удаление & 0,3977 & 0,3998 ± 0,0086 & −0,53 & 1,9220 & 1,9273 ± 0,0107 & −0,28 \\
Очищение & 0,1646 & 0,1642 ± 0,0024 & 0,24 & 1,3117 & 1,3110 ± 0,0036 & 0,05 \\
RCS, новая длительность & 1,0627 & 1,0465 ± 0,0271 & 1,55 & 3,0547 & 3,0361 ± 0,0311 & 0,61 \\
Все каналы, новая длительность & 0,7712 & 0,7722 ± 0,0119 & −0,13 & 2,6625 & 2,6622 ± 0,0109 & 0,01 \\
\bottomrule
\end{tabular}
\caption{Средние при $\delta=0{,}3$. Удаление и очищение рассчитаны методом Такахаси--Таками, повторы --- совместной фазовой моделью. ДИ --- полуширина 95\%-го доверительного интервала по восьми прогонам. «Все каналы» означает одновременный повтор на занятых каналах.}
\label{tab:results}
\end{table}

% tab:results_fixed

\begin{table}[!htbp]
\centering\small\setlength{\tabcolsep}{2pt}
\begin{tabular}{@{}p{4.16cm}p{1.61cm}p{2.83cm}p{1.28cm}p{1.61cm}p{2.83cm}p{1.28cm}@{}}
\toprule
 & \multicolumn{3}{c}{$w_1$} & \multicolumn{3}{c}{$v_1$} \\
Вариант & Расчёт & ИМ $\pm$ ДИ & $\Delta$, \% & Расчёт & ИМ $\pm$ ДИ & $\Delta$, \% \\
\midrule
\multicolumn{7}{l}{Коэффициент вариации: 0,8} \\
RCS, сохранённая длительность & 1,0388 & 1,0364 ± 0,0027 & 0,23 & 3,1695 & 3,1668 ± 0,0029 & 0,09 \\
Все каналы, сохранённая длительность & 1,2726 & 1,2727 ± 0,0036 & −0,01 & 3,4473 & 3,4475 ± 0,0038 & −0,01 \\
\multicolumn{7}{l}{Коэффициент вариации: 1,2} \\
RCS, сохранённая длительность & 1,5888 & 1,5903 ± 0,0053 & −0,09 & 3,7344 & 3,7359 ± 0,0056 & −0,04 \\
Все каналы, сохранённая длительность & 2,2136 & 2,2110 ± 0,0096 & 0,12 & 4,4274 & 4,4245 ± 0,0101 & 0,07 \\
\bottomrule
\end{tabular}
\caption{Средние при повторе с сохранённой длительностью, $\delta=0{,}02$; расчёт с поправкой квадратуры. ДИ --- по 32 прогонам длиной $2\cdot10^6$ уходов.}
\label{tab:results_fixed}
\end{table}

% tab:sd

\begin{table}[!htbp]
\centering\small\setlength{\tabcolsep}{2pt}
\begin{tabular}{@{}p{4.16cm}p{1.61cm}p{2.83cm}p{1.28cm}p{1.61cm}p{2.83cm}p{1.28cm}@{}}
\toprule
 & \multicolumn{3}{c}{$\sigma_W$} & \multicolumn{3}{c}{$\sigma_V$} \\
Вариант & Расчёт & ИМ $\pm$ ДИ & $\Delta$, \% & Расчёт & ИМ $\pm$ ДИ & $\Delta$, \% \\
\midrule
\multicolumn{7}{l}{Коэффициент вариации: 0,8} \\
RCS, удаление & 0,9131 & 0,9340 ± 0,0230 & −2,24 & 1,6980 & 1,7093 ± 0,0157 & −0,66 \\
Очищение & 0,5523 & 0,5555 ± 0,0079 & −0,58 & 1,2655 & 1,2652 ± 0,0052 & 0,02 \\
RCS, новая длительность & 1,9618 & 1,9601 ± 0,0489 & 0,09 & 2,6671 & 2,6657 ± 0,0347 & 0,05 \\
Все каналы, новая длительность & 2,7947 & 2,7967 ± 0,0779 & −0,07 & 3,4332 & 3,4350 ± 0,0638 & −0,05 \\
RCS, сохранённая длительность & 1,7606 & 1,7570 ± 0,0051 & 0,21 & 2,4855 & 2,4824 ± 0,0038 & 0,13 \\
Все каналы, сохранённая длительность & 2,1564 & 2,1554 ± 0,0078 & 0,05 & 2,8566 & 2,8560 ± 0,0064 & 0,02 \\
\multicolumn{7}{l}{Коэффициент вариации: 1,2} \\
RCS, удаление & 1,0268 & 1,0252 ± 0,0268 & 0,16 & 2,0866 & 2,0891 ± 0,0159 & −0,12 \\
Очищение & 0,5699 & 0,5696 ± 0,0086 & 0,05 & 1,4337 & 1,4333 ± 0,0040 & 0,03 \\
RCS, новая длительность & 2,0208 & 1,9836 ± 0,0444 & 1,87 & 3,0849 & 3,0572 ± 0,0319 & 0,91 \\
Все каналы, новая длительность & 1,5746 & 1,5817 ± 0,0274 & −0,45 & 2,6559 & 2,6569 ± 0,0181 & −0,04 \\
RCS, сохранённая длительность & 2,7729 & 2,7800 ± 0,0118 & −0,25 & 3,8517 & 3,8577 ± 0,0094 & −0,15 \\
Все каналы, сохранённая длительность & 3,8815 & 3,8676 ± 0,0195 & 0,36 & 4,8651 & 4,8536 ± 0,0161 & 0,24 \\
\bottomrule
\end{tabular}
\caption{Стандартные отклонения ожидания и пребывания. Варианты и серии ИМ те же, что в таблицах~\ref{tab:results} и~\ref{tab:results_fixed}.}
\label{tab:sd}
\end{table}

% tab:law

\begin{table}[!htbp]
\centering\small\setlength{\tabcolsep}{3pt}
\begin{tabular}{@{}p{1.0cm}p{3.2cm}p{1.6cm}p{1.6cm}p{1.3cm}p{1.6cm}p{1.6cm}p{1.3cm}@{}}
\toprule
 & & \multicolumn{3}{c}{$w_1$} & \multicolumn{3}{c}{$v_1$} \\
$\mathrm{CV}$ & Вариант & $H_2$ & Гамма & $\Delta$, \% & $H_2$ & Гамма & $\Delta$, \% \\
\midrule
0,8 & RCS, удаление & 0,4104 & 0,4100 & 0,11 & 2,0850 & 2,0847 & 0,02 \\
0,8 & Очищение & 0,1925 & 0,1921 & 0,22 & 1,4830 & 1,4831 & 0,00 \\
0,6 & RCS, удаление & 0,4014 & 0,4023 & −0,22 & 2,1362 & 2,1369 & −0,03 \\
0,6 & Очищение & 0,2011 & 0,2021 & −0,48 & 1,5586 & 1,5586 & 0,00 \\
0,5 & RCS, удаление & 0,3947 & 0,3966 & −0,49 & 2,1544 & 2,1559 & −0,07 \\
0,5 & Очищение & 0,2039 & 0,2062 & −1,10 & 1,5912 & 1,5913 & −0,01 \\
\bottomrule
\end{tabular}
\caption{Погрешность замены гамма-распределения законом \(H_2\) с теми же тремя моментами: \(n=3\), \(\rho=0{,}7\), \(\delta=0{,}3\). Столбец «\(H_2\)» --- метод Такахаси--Таками, «Гамма» --- PH-представление с 16 узлами квадратуры. \(\Delta=100(H_2-\text{гамма})/\text{гамма}\).}
\label{tab:law}
\end{table}
